\chapter{2005年重庆卷（理）}
\section{单选题}
\begin{problem}
因 \((x + 2)^{2} + y^{2} = 5\) 关于原点(0,0)对称的圆的方程为（\quad）

\choices
    {\((x - 2)^{2} + y^{2} = 5\)}
    {\(x^{2} + (y - 2)^{2} = 5\)}
    {\((x + 2)^{2} + (y + 2)^{2} = 5\)}
    {\(x^{2} + (y + 2)^{2} = 5\)}
\end{problem}

\begin{problem}
\(\left(\frac{1 + 1}{1 - 1}\right)^{2005} =\)（\quad）

\choices
    {1}
    {\(-1\)}
    {\(2^{2005}\)}
    {\(-2^{2005}\)}
\end{problem}

\begin{problem}
若函数 \( f(x) \) 是定义在 \( \R \) 上的偶函数，在 \((- \infty, 0]\) 上是减函数，且 \( f(2) = 0 \)，则使得 \( f(x) < 0 \) 的 \( x \) 的取值范围是（\quad）

\choices
    {\((-\infty, 2)\)}
    {\((2, +\infty)\)}
    {\((-\infty, -2) \cup (2, +\infty)\)}
    {\((-2, 2)\)}
\end{problem}

\begin{problem}
已知 \( A(3,1) \) ， \( B(6,1) \) ， \( C(4,3) \) ，D 为线段 BC 的中点，则向量 \( \overrightarrow{AC} \) 与 \( \overrightarrow{BD} \) 的夹角为（\quad）

\choices
    {\(\frac{5}{4} -\arccos \frac{4}{5}\)}
    {\(\arccos \frac{4}{5}\)}
    {\(\arccos \left(-\frac{4}{5}\right)\)}
    {\(-\arccos \left(-\frac{4}{5}\right)\)}
\end{problem}

\begin{problem}
若 \(x, y\) 是正数，则 \(\left(x + \frac{1}{2y}\right)^2 + \left(y + \frac{1}{2x}\right)^2\) 的最小值是（\quad）

\choices
    {3}
    {\( \frac{7}{2} \)}
    {4}
    {\(\frac{9}{2}\)}
\end{problem}

\begin{problem}
已知 \(\alpha\)、\(\beta\) 均为锐角，若 \(p: \sin \alpha < \sin (\alpha + \beta), q: \alpha + \beta < \frac{\pi}{2}\)，则 \(p\) 是 \(q\) 的（\quad）

\choices
    {充分而不必要条件}
    {必要而不充分条件}
    {充要条件}
    {既不充分也不必要条件}
\end{problem}

\begin{problem}
对于不重合的两个平面 \(\alpha\) 与 \(\lambda\)，给定下列条件: \circled{1} 存在平面 \(\gamma\)，使得 \(\alpha\)、\(\beta\) 都垂直于 \(\gamma\)； \circled{2} 存在平面 \(\gamma\)，使得 \(\alpha\)、\(\beta\) 都平行于 \(\gamma\)； \circled{3} \(\alpha\) 内有不共线的三点到 \(\beta\) 的距离相等； \circled{4} 存在异面直线 \(l\)，\(m\)，使得 \(l \nmid \alpha, l \nmid \beta, m \nmid \alpha, m \nmid \beta\). 其中，可以判定 \(\alpha\) 与 \(\beta\) 平行的条件有（\quad）

\choices
    {1 个}
    {2 个}
    {3 个}
    {4 个}
\end{problem}

\begin{problem}
若 \(\left(2x - \frac{1}{x}\right)^n\) 展开式中含 \(\frac{1}{x^2}\) 项的系数与含 \(\frac{1}{x^4}\) 项的系数之比为\(-5\)，则 \(n\) 等于（\quad）

\choices
    {4}
    {6}
    {8}
    {10}
\end{problem}

\begin{problem}
若动点 \((x, y)\) 在曲线 \(\frac{x^2}{4} + \frac{y^2}{b^2} = 1 (b > 0)\) 上变化，则 \(x^2 + 2y\) 的最大值为（\quad）

\choices
    {\(\left\{ \begin{array}{ll} \frac{b^2}{4} + 4, & 0 < b < 4, \\ 2b, & b \geqslant 4. \end{array} \right.\)}
    {\(\left\{ \begin{array}{ll} \frac{b^2}{4} + 4, & 0 < b < 2, \\ 2b, & b \geqslant 2. \end{array} \right.\)}
    {\( \frac{b^{2}}{4} + 4 \)}
    {2b}
\end{problem}

\begin{problem}
在体积为 1 的三棱锥 A-BCD 侧棱 AB、AC、AD 上分别取点 E、F、G，使 AE:EB = AF:FC = AG:GD = 2:1，记 O 为三平面 BCG、CDE、DBF 的交点，则三棱锥 O-BCD 的体积等于（\quad）

\choices
    {\( \frac{1}{9} \)}
    {\( \frac{1}{8} \)}
    {\( \frac{1}{7} \)}
    {\( \frac{1}{4} \)}
\end{problem}

\section{填空题}
\begin{problem}
已知集合 \(A = \{x \in \R \mid x^2 - x - 6 < 0\}\)，集合 \(B = \{x \in \R \mid |x - 2| < 2\}\)，则 \(A \cap B =\)  \fillinblank{}.
\end{problem}

\begin{problem}
曲线 \( y = x^{2} \) 在点 \( (a, a^{4}) \) （\( a \neq 0 \)）处的切线与 \( x \) 轴、直线 \( x = a \) 所围成的三角形的面积为 \( \frac{1}{4} \)，则 \( a = \)  \fillinblank{}.
\end{problem}

\begin{problem}
已知 \(\alpha\)、\(\beta\) 均为锐角，且 \(\cos (\alpha + \beta) = \sin (\alpha - \beta)\)，则 \(\tan \alpha = \)  \fillinblank{}.
\end{problem}

\begin{problem}
\(\lim_{n\to \infty}\frac{2^{2n} - 3^{2n + 1}}{2^{2n} + 3^{2n}} =\)
\end{problem}

\begin{problem}
某轻轨列车有 4 节车厢，现有 6 位乘客准备乘坐，设每一位乘客进入每节车厢是等可能的，则这 6 位乘客进入各节车厢的人数恰好为 0, 1, 2, 3 的概率为 \fillinblank{} \fillinblank{}.
\end{problem}

\begin{problem}
连接抛物线上任意四点组成的四边形可能是 \fillinblank{}. (填写所有正确选项的序号) \circled{1}菱形；\circled{2}有3条边相等的四边形；\circled{3}梯形；\circled{4}平行四边形；\circled{5}有一组对角相等的四边形 \fillinblank{}.
\end{problem}

\section{解答题}
\begin{problem}
若函数 \( f(x) = \frac{1 + \cos 2x}{4\sin\left(\frac{\pi}{2} + x\right)} - a\sin \frac{\pi}{2}\cos \left(\pi - \frac{\pi}{2}\right) \) 的最大值为2，试确定常数 \( a \) 的值.
\end{problem}

\begin{problem}
在一次购物抽奖活动中，假设某 10 张券中有一等奖券 1 张，可获价值 50 元的奖品; 有一二等奖券 3 张. 每张可获价值 10 元的奖品; 其余 6 张没有奖，某顾客从距 10 张券中任抽 2 张，求: (1)该顾客中奖的概率; (2) 该顾客获得的奖品总价值 \(\xi\) (元) 的概率分布列和期望 \(E\xi\).
\end{problem}

\begin{problem}
已知 \(a \in \R\)，讨论函数 \(f(x) = \e^{x}(x^{2} + ax + a + 1)\) 的极值点的个数.
\end{problem}

\begin{problem}
如图，在三棱柱 \(ABC - A_{1}B_{1}C_{1}\) 中，\(AB \perp\) 侧面 \(BB_{1}C_{1}C\)，\(E\) 为棱 \(CC_{1}\) 上异于 \(C\)，\(C_{1}\) 的一点，\(EA \perp EB_{1}\)，已知 \(AB = \sqrt{2}\)，\(BB_{1} = 2\)，\(BC = 1\)，\(\angle BCC_{1} = \frac{\pi}{3}\)，求: (1) 异面直线 AB 与 \( EB_{1} \) 的距离; (2) 二面角 \(A - EB_{1} - A_{1}\) 的平面的角的正切值.
\end{problem}

\begin{problem}
已知椭圆 \(C_1\) 的方程为 \(\frac{x^2}{4} + y^2 = 1\)，双曲线 \(C_2\) 的左、右焦点分别为 \(C_1\) 的左、右顶点，而 \(C_2\) 的左、右顶点分别是 \(C_1\) 的左、右焦点. （1）求双曲线 \(C_2\) 的方程. (2) 若直线 \(l: y = kx + \sqrt{2}\) 与椭圆 \(C_1\) 及双曲线 \(C_2\) 都恒有两个不同的交点，且 \(l\) 与 \(C_2\) 的两个交点 \(A\) 和 \(B\) 满足 \(\overrightarrow{OA} \cdot \overrightarrow{OB} < 6\) （其中 \(O\) 为原点），求 \(k\) 的取值范围.
\end{problem}

\begin{problem}
数列 \(\{a_{n}\}\) 满足 \(a_1 = 1\) 且 \(a_{n + 1} = \left(1 + \frac{1}{n^2 + n}\right)a_n + \frac{1}{2^n}\) (\(n \geqslant 1\)). (1) 用数学归纳法证明： \( a_{n} \geqslant 2 (n \geqslant 2) \) ; (2) 已知不等式 \(\ln(1+x) 0\) 成立，证明 \( a_{n}<\e^{x} \) ( \( n\geqslant1 \) )，其中无理数是 \( \e=2718288\cdots \) .
\end{problem}

